\documentclass[12pt]{amsart}
\usepackage[linkcolor=blue, citecolor = blue]{hyperref}
\hypersetup{colorlinks=true}

\usepackage{amsthm,amsfonts,amsmath,amssymb,amscd} 

\usepackage{verbatim}
\usepackage{float}
\usepackage{wrapfig}
\usepackage{mathtools}
\usepackage[color=orange!20, linecolor=orange]{todonotes}
\newcommand{\marginnote}{\todo[color=blue!20, linecolor=blue, bordercolor=blue]}
\newcommand{\TODO}{\todo[inline]}
\usepackage{caption, subcaption}
\usepackage{enumitem}

\usepackage{ytableau}

%\usepackage{tikz}
%\usetikzlibrary{decorations.markings}
%\tikzstyle{vertex}=[circle, draw, inner sep=0pt, minimum size=4pt]
%\newcommand{\vertex}{\node[vertex]}
%\newcounter{Angle}
%\usetikzlibrary{arrows,automata}
%\usepackage{tikz, tikz-3dplot, pgfplots}
%\usepackage{tkz-graph}
%\usetikzlibrary[positioning,patterns]
%\usepackage{tikz-3dplot}
%\usepackage{youngtab}
%\usepackage{genyoungtabtikz}
%\usepackage{gfsdidot}
%\usepackage{concrete}
\usepackage{wasysym}

\newtheorem{theorem}{Theorem}[section]
\newtheorem{proposition}[theorem]{Proposition}
\newtheorem{lemma}[theorem]{Lemma}
\newtheorem{corollary}[theorem]{Corollary}
 
\theoremstyle{definition}
\newtheorem{definition}[theorem]{Definition}
\newtheorem{example}[theorem]{Example}
\newtheorem{problem}[theorem]{Problem}
\newtheorem{conjecture}[theorem]{Conjecture}

\theoremstyle{remark}
\newtheorem{remark}{Remark}

\newcommand{\sgn}{\mathrm{sgn}}

\setlength{\textwidth}{6.5in}
\setlength{\oddsidemargin}{0.0in}
\setlength{\evensidemargin}{0.0in}
\setlength{\topmargin}{-0.0in}
\setlength{\textheight}{8.5in}
\renewcommand{\baselinestretch}{1.1}


%%%%%%%%%%%%%%%%%%%%%%
\title[Bounds on the number of higher-dimensional partitions]{Bounds on the number of higher-dimensional partitions} 

\author[Damir Yeliussizov]{Damir Yeliussizov}
%\thanks{\today}
%\thanks{* Department ???. Email: \texttt{yeldamir@gmail.com}}

\address{KBTU, Almaty, Kazakhstan}
\email{\href{mailto:yeldamir@gmail.com}{yeldamir@gmail.com}}
%\email{\href{mailto:yeldamir@gmail.com}{yeldamir@gmail.com}}
%\keywords{Higher-dimensional partitions, MacMahon's numbers}
\subjclass[2020]{05A16, 05A17, 05A20, 11P21}

%\date{June, 2022}

\begin{document}

\begin{abstract}
We establish some bounds on the number of higher-dimensional partitions by volume. 
In particular, we give bounds via vector partitions and MacMahon's numbers. %comparison with MacMahon's numbers.
\end{abstract}

\maketitle

%\tableofcontents

\section{Introduction}

{Higher-dimensional partitions} were introduced by P. A. MacMahon \cite{macmahon} as a natural generalization of  
(one-dimensional) 
integer partitions and (two-dimensional) 
plane partitions. %, due to. 
While these  objects %are now known to 
arise in different 
fields such as %including 
algebra, geometry, combinatorics,  and statistical physics 
(%various aspects %of higher-dimensional partitions 
%can be found in 
see e.g. 
\cite{mr, bgp, bbs, gov, ms, nekrasov} on some aspects), %of higher-dimensional partitions. 
%On the other hand, 
not much is known about 
higher-dimensional partitions 
%them 
in dimensions three or more (cf. \cite[\textsection 7.20]{sta}), including %In particular, %one of the open problems concerns 
their asymptotics. %behavior 


\vspace{0.5em}

In this note, we prove some asymptotic lower and upper bounds on the number of $d$-dimensional partitions of volume $n$.

%\vspace{0.5em}

\subsection{Definition} %Higher-dimensional partitions} 
A {\it $d$-dimensional partition} is a hypermatrix 
$\pi = (\pi_{i_1,\ldots, i_d})_{i_1,\ldots, i_d \ge 1}$ of nonnegative integers with only finitely many nonzero entries such that %the following inequalities hold:
$
\pi_{i_1,\ldots, i_d} \ge \pi_{j_1,\ldots, j_d} \text{ for all } i_1 \le j_1, \ldots, i_d \le j_d.
$
The {\it volume} (or {\it size}) of $\pi$ is $|\pi| = \sum_{i_1,\ldots, i_d} \pi_{i_1,\ldots, i_d}.$ 

\vspace{0.5em}

For $d = 1$, this definition corresponds to the usual integer partitions. For $d = 2$, it defines %the known combinatorial object of 
{plane partitions}. We refer e.g. to \cite{sta2, andrews, sta, krat} on the theory of these well-studied objects.

\vspace{0.5em}

Let $\mathsf{p}_d(n)$ be the number of $d$-dimensional partitions of volume $n$. 

\subsection{MacMahon's conjecture and asymptotics} 
%Higher-dimensional partitions were introduced by MacMahon \cite{mac}. Consider the generating function.
{\it MacMahon's numbers} $\mathsf{m}_d(n)$ are defined via the generating series
$$
\sum_{n \ge 0} \mathsf{m}_d(n)\, x^n = \prod_{n \ge 1} (1 - x^{n})^{-\binom{n + d - 2}{d - 1}}.
$$
MacMahon \cite{macmahon} conjectured that $\mathsf{p}_d(n) = \mathsf{m}_d(n)$ which is true for $d = 1,2$, but turns out to be false for $d \ge 3$ as shown %by Atkin, Bratley, Macdonald and McKay 
in \cite{atkin}; 
in fact, %one can actually compute that 
$
\mathsf{m}_d(6) - \mathsf{p}_d(6) = \binom{d}{3} + \binom{d}{4}
$
(see \cite{andrews}).

%\vspace{0.5em}

An interpretation of MacMahon's numbers via $d$-dimensional partitions was given in \cite{ay}, where it is shown that (instead of the volume) $\mathsf{m}_d(n)$ computes %the number of 
$d$-dimensional partitions by another statistic called the {\it corner-hook volume}, which we recall in  \textsection\ref{secmac}.

%\subsection{Conjectured asymptotics} 
\vspace{0.5em}
%\newpage

Throughout the paper, we consider the asymptotics as $n \to \infty$ assuming $d$ is fixed.

It is known that (see \cite{bgp})
$$
\lim_{n \to \infty} \frac{\log \mathsf{m}_d(n)}{n^{d/(d  + 1)}} = \gamma_d = {{d^{-d/(d+1)} (d+1)\, \zeta(d+1)^{1/(d + 1)}}},
$$
where $\zeta(s) = \sum_{n \ge 1} n^{-s}$ is the Riemann zeta function. (Using the explicit generating function, a more detailed asymptotics of $\mathsf{m}_d(n)$ can be found.) 

Concerning the asymptotics of $\mathsf{p}_d(n)$, 
it is %also 
known from 
\cite{bhatia} that $\log \mathsf{p}_{d}(n) = \Theta(n^{d/(d+1)})$. 
(Note that no explicit generating function is known for $\mathsf{p}_d(n)$ when $d \ge 3$.) 

Even though MacMahon's conjecture is false, based on some numerical simulations %experiments 
for (3-dimensional) solid partitions, 
it was conjectured in \cite{mr} 
%conjectured 
that it might hold asymptotically, i.e. $\log \mathsf{p}_3(n) \sim \log \mathsf{m}_3(n)$. 
Based on experiments extended to higher dimensions, it was conjectured more generally in \cite{bgp} that
$\log \mathsf{p}_d(n) \sim \log \mathsf{m}_d(n) \sim \gamma_d\, n^{d/(d + 1)}.$
%proposed the following more general conjecture. 
However, later experiments for solid partitions reported in \cite{dg} suggested that the conjecture of \cite{mr} appears to be false and $\mathsf{p}_3(n)$ grows faster than $\mathsf{m}_3(n)$.


%Note that in contrast with MacMahon's numbers, 
%Note that no formulas or explicit generating functions are known for $\mathsf{p}_d(n)$ when $d \ge 3$. %, which makes it difficult to understand its asymptotics. 
%The difficulty in understanding the asymptotics of %the sequence $\mathsf{p}_d(n)$ is mainly in the absence of any formulas or explicit generating functions for it. %, which is how usually asymptotics would be found. 

\subsection{Asymptotic bounds} 


%Let $p_d(n)$ be the number of $d$-dimensional partitions of size (volume) $n$ (see definitions in \textsection \ref{?}). 

%\subsection{Asymptotic bounds} 
%\vspace{0.5em}

%More explicitly, 
We prove the following asymptotic bounds.
%We prove 

\begin{theorem}\label{thm:one}
%Let $d > 1$. 
%As $n \to \infty$,
For all sufficiently large $n$, 
we have the following bounds: 
$$
\alpha_d \, {n^{d/(d+1)}} \le {\log \mathsf{p}_d(n)}  \le \beta_d \, {n^{d/(d+1)}} + d \log n, %(d+1) \zeta(d+1)^{1/(d + 1)} 
%+ o(1),
$$
where the constants are
$$
\alpha_d = \frac{(d+1)}{(d+1)!^{1/(d+1)}}  %\left(\log 2 + {3^{-d}} \log \frac{3}{2} \right), %
\cdot \delta_d ~\text{ for some } \delta_d > \log 2, 
\qquad \beta_d = (d+1) \zeta(d+1)^{1/(d + 1)}. %= d^{d/(d+1)} \gamma_d.
$$
\end{theorem}

This theorem %Theorem~\ref{thm:one} 
improves an asymptotic result of \cite{bhatia} where it was shown that $\log \mathsf{p}_{d}(n) = \Theta(n^{d/(d+1)})$ without explicit bounds (which can be computed to give much wider range for the constants). 
%\begin{remark}

After an earlier version of 
this paper was written, we learned about the very recent work \cite{dai},
where a slightly smaller lower bound constant $\alpha'_d = \frac{(d+1)}{(d+1)!^{1/(d+1)}} \log 2 - o(1)$ and 
a significantly larger upper bound  $\gamma_1 (d+1)^{\log(d+1)}$ %\pi \sqrt{2/3}
(which optimizes the bound from \cite{bhatia}) were obtained. %\end{remark}
%Let us note that 
%our approach in proving 
For the upper bound %is based on a comparison with vector partitions (see \textsection\ref{secvec}) which 
our approach
is different from \cite{bhatia, dai}; for the lower bound 
our estimate %it
is similar to \cite{dai}, but the technique is different giving the constant $\alpha_d$ (cf. Theorem~\ref{coral}).

Comparing the lower bound with MacMahon's numbers, we have $\alpha_d > \gamma_d$ for $d \ge 7$  which %our lower bound $\alpha_d$ 
shows %Conjecture~\ref{conj1} 
%and %shows 
that {$\mathsf{p}_d(n)$ is asymptotically larger than $\mathsf{m}_d(n)$ in dimensions $d \ge 7$}, i.e. %, as stated in the following result.
%\begin{corollary}\label{thm1}
%For $d \ge 7$ we have
$$%\begin{align}\label{thm1}
\liminf_{n \to \infty} 
\frac{\log \mathsf{p}_{d}(n)}{n^{d/(d+1) }} > \gamma_d %\eqno{(\diamond)}
%\lim_{n \to \infty} 
%\frac{\log \mathsf{m}_{d}(n)}{n^{d/(d+1) }}.
$$%\end{align}
and hence the conjecture of \cite{bgp} is false.
We note that this was shown in \cite{dai} using the constant $\alpha'_d$. %> \gamma_d$)
In dimensions $3 \le d \le 6$, an asymptotic comparison of $\mathsf{p}_d(n)$ and $\mathsf{m}_d(n)$ remains unknown. 


\begin{remark}
Interestingly, all available numerical values of $\mathsf{p}_d(n)$ (see \cite{atkin, knuth, boltzmann, ekhad} and \cite[A000293, A000294]{oeis}) show that $\mathsf{m}_d(n) \ge \mathsf{p}_d(n)$ and that the sequences grow relatively close. The above limit inequality %~($\diamond$) %\eqref{thm1} 
tells that this is delusive for asymptotics. 
%In fact, later experiments for solid partitions reported in \cite{dg} suggested that the conjecture of \cite{mr} can be false and $\log \mathsf{p}_3(n) \gg \log \mathsf{m}_3(n)$, and so ($\diamond$) %\eqref{thm1} 
%agrees with this in higher dimensions.
\end{remark}


%\vspace{0.5em}
%To be more precise, we can take the constant $\delta_d$ in Theorem~\ref{thm:one} to be e.g. (cf. Corollary~\ref{coral})
%$$
%\delta_d \ge \log 2 + {3^{-d}} \log \frac{3}{2} %\frac{1}{2^d} \log \left(2^{2^d} + 2^{2^d - d - 1} \right) %\log 2 + \frac{1}{2^d} \log \frac{3}{2} %\log \left(2^{2^d}  +  \sum_{i = 0}^{d} 3^{\binom{d}{i}} - 2^{\binom{d}{i }} \right) 
%- \varepsilon \text{ for any } \varepsilon > 0.
%$$

\begin{remark}
It is also open to prove the existence of 
$\lim_{n \to \infty } {n^{-d/(d+1)}} {\log \mathsf{p}_d(n)}$
(which is assumed in some literature referring to \cite{bhatia}). 
\end{remark}

\subsection{An upper bound via vector partitions} 
 
To prove the above results (specifically an upper bound) we give a comparison with numbers whose explicit generating functions are known. %obtain some combinatorial bounds. %shown below. 

%Let $a_d(k)$ be the number of $d$-dimensional partitions whose diagram (see definition in the next section) lies inside the simplex $x_1 + \ldots + x_{d+1} \le k$ in $\mathbb{Z}_+^{d+1}$. Then the following basic bounds hold: %It can be shown that (see \textsection\ref{?}) the following basic bounds hold:
%$$
%2^{\binom{k-1}{d}} \le a_d(k) \le 2^{2 \binom{k-1}{d}},
%$$
%see \textsection\ref{secsimpl} where we also show an improved lower bound. 

{\it Vector partition} (or {\it multipartition}) {\it numbers} $\mathsf{p}(n_1,\ldots, n_d)$ are defined via the multivariate generating series 
$$
\sum_{n_1,\ldots, n_d } \mathsf{p}(n_1,\ldots, n_d)\, x_1^{n_1} \cdots x_d^{n_d}  = \prod_{i_1,\ldots, i_d \ge 1} \left(1 - x_1^{i_1} \cdots x_d^{i_d} \right)^{-1}.
$$
(Note that for $d = 1$, we have $\mathsf{p}(n) = \mathsf{p}_1(n)$ is the usual number of partitions.) 


\begin{theorem}\label{thmc}
The following inequality holds:
$$
\mathsf{p}_d(n) \le n^{d}\, {\mathsf{p}}(\underbrace{n, \ldots, n}_{d \text{ times}}). %, \qquad \mathsf{p}_d(n) < dn \cdot \mathsf{m}_d(dn) %\le? dn\cdot m_d(dn).
$$
\end{theorem}

To obtain this
upper bound %is more interesting. 
%In particular, %to prove the upper bound, 
we give a new combinatorial interpretation of vector partition numbers via $d$-dimensional partitions (see Lemma~\ref{wp}).
We also obtain an upper bound via MacMahon's numbers: $$\mathsf{p}_d(n) < dn \cdot \mathsf{m}_d(dn).$$


Our proofs %of Theorem~\ref{thmc} 
use some tools developed in \cite{ay}.  


\section{Preliminaries}
%\subsection{Partitions, shapes, and diagrams}

We use the following basic notation: $\mathbb{N}$ is the set of nonnegative integers; $\mathbb{Z}_+$ is the set of positive integers; $\{\mathbf{e}_1, \ldots, \mathbf{e}_d\}$ is the standard basis of $\mathbb{Z}^d$; and $[n] := \{1, \ldots, n \}$.

\vspace{0.5em}

%\subsection{Partitions and diagrams} 
A {\it $d$-dimensional $\mathbb{N}$-hypermatrix} is an array $\left(a_{i_1, \ldots, i_d}\right)_{i_1, \ldots, i_d \ge 1}$ of nonnegative integers %$a_{\mathbf{i}}$ 
with finite support (i.e. with only finitely many nonzero elements). 
A {\it $d$-dimensional partition} is a $d$-dimensional 
%hypermatrix 
$\mathbb{N}$-hypermatrix 
$\left(\pi_{i_1, \ldots, i_d} \right)_{i_1,\ldots, i_d \ge 1}$ 
%of nonnegative integers
%with finite support (i.e. only finitely many nonzero elements),
such that 
$$
\pi_{i_1, \ldots, i_d} \ge \pi_{j_1, \ldots, j_d} ~\text{ for }~ i_1 \le j_1, \ldots, i_d \le j_d.
$$
Let %$\mathcal{M}^{(d)}$ be the set of $d$-dimensional $\mathbb{N}$-hypermatrices and 
$\mathcal{P}^{(d)}$ be the set of $d$-dimensional partitions. 
For $\pi = (\pi_{i_1, \ldots, i_d}) \in \mathcal{P}^{(d)}$, the {\it volume} (or size) of $\pi$ denoted by $|\pi|$ is defined as 
$$
|\pi| = \sum_{i_1, \ldots, i_d} \pi_{i_1, \ldots, i_d}.
$$

Any partition $\pi$ is uniquely determined by its {\it diagram} $D(\pi)$ which is the set 
$$D(\pi) := \{(i_1, \ldots, i_d, i) \in \mathbb{Z}^{d+1}_+: 1 \le i \le \pi_{i_1, \ldots, i_d} \} %\subset \mathbb{Z}^{d+1}_+.
$$
of cardinality $|D(\pi)| = |\pi|$. 

Let us note that 
for a set $\rho \subset \mathbb{Z}^{d+1}_+$, the following conditions are equivalent:
\begin{itemize}
\item[(1)] The set $\rho$ is the diagram of some $d$-dimensional partition.
\item[(2)] The set $\rho$ is finite and has the property that if $\mathbf{i} \in \mathbb{Z}^{d+1}_+$  and $\ell \in [d+1]$ satisfy $\mathbf{i} + \mathbf{e}_{\ell} \in \rho$, then $\mathbf{i} \in \rho$. 
\end{itemize}
Sometimes we may identify a partition with its diagram. 

\vspace{0.5em}

We denote by $\mathsf{p}_d(n)$ the number of $d$-dimensional partitions of size $n$. Let also $\widetilde{\mathsf{p}}_d(n)$ be the number of $d$-dimensional partitions of size at most $n$, so that for $ n > 1$ we have
$$
\mathsf{p}_d(n) < \widetilde{\mathsf{p}}_d(n) = \sum_{i \le n}^{} \mathsf{p}_d(i) \le n\cdot \mathsf{p}_d(n),
$$
and hence $\log \widetilde{\mathsf{p}}_d(n) = \log \mathsf{p}_d(n) + O(\log n)$ 
has the same main term of asymptotics as $n \to \infty$; in our case we have more precisely
$$\frac{\log \widetilde{\mathsf{p}}_{d}(n)}{n^{d/(d+1)}} = \frac{\log \mathsf{p}_{d}(n)}{n^{d/(d+1)}} + o(1).$$ 

\section{A map from hypermatrices to partitions}\label{secphi}

In this section we describe a map from hypermatrices to partitions and some of its properties studied in \cite{ay} which will be useful for proving our results.

\vspace{0.5em}

A lattice path in $\mathbb{Z}^d$ is called {\it directed} if it uses only steps of the form $\mathbf{i} \to \mathbf{i} + \mathbf{e}_{\ell}$ for $\mathbf{i} \in \mathbb{Z}^d$ and $\ell \in [d]$. 
%Let $\mathcal{M}^{(d)}$ be the set 
For a $d$-dimensional $\mathbb{N}$-hypermatrix $A = (a_{i_1, \ldots, i_d})_{i_1, \ldots, i_d \ge 1}$ we define the following map: 
$$\Phi : A \longmapsto \pi = (\pi_{i_1,\ldots, i_d}), \qquad \pi_{i_1,\ldots, i_d} = \max_{\Pi : (i_1,\ldots, i_d) \to (\infty^d)} \sum_{(j_1,\ldots, j_d)\in \Pi } a_{j_1,\ldots, j_d},$$
where the sum is over directed lattice paths $\Pi$ which start at $(i_1,\ldots, i_d)$. Equivalently, we also have 
$$
\pi_{\mathbf{i}} = a_{\mathbf{i}} + \max_{\ell \in [d]} \pi_{\mathbf{i} + \mathbf{e}_{\ell}}, \quad \mathbf{i} \in \mathbb{Z}_+^{d}.
$$
Note that 
%if $A  = (a_{i_1, \ldots, i_d})$ is a hypermatrix of nonnegative integers with finite support (i.e. with only finitely many nonzero elements), then 
$\pi = \Phi(A) \in \mathcal{P}^{(d)}$ is a $d$-dimensional partition.

\vspace{0.5em}

%This map has the following important property.
%\begin{theorem}[\cite{ay}] 
%Importantly, the map $\Phi$ 
%(i);
%(ii) %The map $\Phi$ 
%is a bijection between $d$-dimensional $\mathbb{N}$-hypermatrices %of nonnegative integers with finite support %$\mathcal{M}^{(d)}$ 
%and $d$-dimensional partitions \cite{ay}. %$\mathcal{P}^{(d)}$.
%In particular, it is injective on a set of $\mathbb{N}$-hypermatrices.
%(ii) Let $\pi \in \mathcal{P}^{(d)}$ and $A = \Phi^{-1}(\pi)$. Then $a_{\mathbf{i}} = |\{(\mathbf{i}, i_{d+1}) : (\mathbf{i}, i_{d+1}) \in \mathrm{Cor}(\pi) \}|$
%\end{theorem}
%In particular, it is injective on any set of $\mathbb{N}$-hypermatrices. 
%Corners???

%\subsection{Corners}
For a partition $\pi \in \mathcal{P}^{(d)}$, define the set of {\it corners} of $\pi$ as follows:
$$
\mathrm{Cor}(\pi) := \left\{\mathbf{i} \in \mathbb{Z}^{d+1}_+ : \mathbf{i} \in D(\pi),\, \mathbf{i} + \mathbf{e}_{\ell} \not\in D(\pi) \text{ for all } \ell \in [d] \right\}.
$$
(Here $\{e_{\ell}\}$ is the standard basis in $\mathbb{Z}^{d+1}$.) For example, if $\pi \in \mathcal{P}^{(2)}$ is a plane partition then $$\mathrm{Cor}(\pi) = \{(i,j,k) : (i,j, k) \in D(\pi), (i+1, j, k), (i, j+1, k) \not\in D(\pi) \}.$$ %Let $\mathrm{cor}(\pi) := |\mathrm{Cor}(\pi)|$ be the number of corners of $\pi$. 

%The inverse map $\Phi^{-1}$ can be described via projection of corners as follows. Let $\pi \in \mathcal{P}^{(d)}$ and $A = (a_{i_1,\ldots, i_d}) = \Phi^{-1}(\pi)$. Then we have 
%$$
%a_{i_1,\ldots, i_d} = \left| \left\{i : (i_1,\ldots, i_d, i) \in \mathrm{Cor}(\pi) \right\} \right|.
%$$

\vspace{0.5em}

The map $\Phi$ has the following important properties.

\begin{theorem}[\cite{ay}] The following properties hold:

(i) The map $\Phi$  is a bijection between $d$-dimensional $\mathbb{N}$-hypermatrices and $d$-dimensional partitions. In particular, it is injective on any set of $\mathbb{N}$-hypermatrices. 

(ii) The inverse map $\Phi^{-1}$ can be described via projection of corners as follows. For $\pi \in \mathcal{P}^{(d)}$ and $A = (a_{i_1,\ldots, i_d}) = \Phi^{-1}(\pi)$, we have 
$$
a_{i_1,\ldots, i_d} = \left| \left\{i : (i_1,\ldots, i_d, i) \in \mathrm{Cor}(\pi) \right\} \right|.
$$
\end{theorem}

As an application, the following multivariate identity was proved in \cite{ay}: %using the map $\Phi$. 
%\begin{theorem}[\cite{ay}]
%The following multivariate identity holds:
$$%\begin{align}\label{id1}
\sum_{\pi \in \mathcal{P}^{(d)}} \prod_{(i_1,\ldots, i_d, i_{d+1}) \in \mathrm{Cor}(\pi)} x_{i_1,\ldots, i_d} = \prod_{i_1,\ldots, i_d \ge 1} \left(1 - x_{i_1,\ldots, i_d} \right)^{-1}. \eqno{(\diamond)}
$$%\end{align}
%\end{theorem}


\section{Bounds on partitions inside a simplex}\label{secsimpl}
Let $\Delta_d(k)$ be the set of positive integer points inside the $(d+1)$-dimensional simplex  $x_1 + \ldots + x_{d+1} \le k$, i.e.
$$
\Delta_d(k) := \left\{(x_1,\ldots, x_{d+1}) \in \mathbb{Z}_{+}^{d+1} :  %(i_1, \ldots, i_{d+1}) \in D(\pi) \implies 
x_1 + \ldots + x_{d+1} \le k \right\}.
$$
Note that $\Delta_d(k)$ is a diagram of a $d$-dimensional partition of size (for $k \in \mathbb{N}$)
$$
%\frac{(k-d)^{d+1}}{(d+1)!} \le 
|\Delta_d(k)| = \binom{k}{d+1} = \frac{k^{d+1}}{(d+1)!} - O(k^d) \le \frac{k^{d+1}}{(d+1)!}.
$$ 

Let 
$$
A_d(k) := \left\{\pi \in \mathcal{P}^{(d)} :  D(\pi) \subseteq \Delta_d(k)
%(i_1, \ldots, i_{d+1}) \in D(\pi) \implies i_1 + \ldots + i_{d+1} \le k 
\right\}
$$
be the set of $d$-dimensional partitions whose diagram lies inside $\Delta_d(k)$ %the simplex $x_1 + \ldots + x_{d+1} \le k$ 
and denote by $a_d(k) := |A_d(k)|$ the number of such partitions.
(For example, $a_1(k)$ is the number of partitions whose diagram lies inside the triangle $\{(x,y) \in \mathbb{Z}^{2}_+ : x + y \le k\}$, and so $a_1(k)$ is the %$k$-th 
Catalan number.)

\vspace{0.5em}

Let $B_d(k)$ be the set of hypermatrices $B = (b_{\mathbf{i}})_{\mathbf{i} \in \Delta_{d - 1}(k - 1)}$ (with support in $\Delta_{d - 1}(k - 1)$) with entries $b_{\mathbf{i}} \in \{0,1,2 \}$ satisfying the following property: if $b_{\mathbf{i}} = 2$ then $\mathbf{i} + \mathbf{e}_{\ell}\in \Delta_{d - 1}(k - 1)$ and $b_{\mathbf{i} + \mathbf{e}_{\ell}} = 0$ for all $\ell \in [d]$. Denote by $b_d(k) := |B_d(k)|$ the number of such hypermatrices. 

\begin{lemma}
We have: 
$$
a_d(k) \ge b_d(k).
$$
\end{lemma}

\begin{proof}
Let us show that $\Phi : B_d(k) \to A_d(k)$. 
For  $B = (b_{i_1,\ldots, i_d}) \in B_d(k)$ let $\pi = (\pi_{i_1,\ldots, i_d}) = \Phi(B)$.  
Note that $\pi \in \mathcal{P}^{(d)}$ is a $d$-dimensional partition. Let us show that $\pi \in A_d(k)$.  
We are going to prove by induction on $m = k - i_1 - \ldots - i_d \ge 0$ that $\pi_{i_1,\ldots, i_d} \le m$. For $m = 0$ the statement is obvious. For $m = 1$ we have $(i_1,\ldots, i_d) + \mathbf{e}_{\ell} \not\in \Delta_{d-1}(k-1)$ for all $\ell \in [d]$ and hence $\pi_{i_1,\ldots, i_d} = b_{i_1,\ldots, i_d} \le 1$. %since $(i_1, \ldots, i_d)$ is either not in any box $J$ (and  $b_{i_1,\ldots, i_d} = 0$) or $(i_1,\ldots, i_d)\ne (2j_1 - 1, \ldots, 2j_d - 1)$ for a box $J$ (and $b_{i_1,\ldots, i_d} \le 1$).  
Assume now the statement holds for all $i'_1,\ldots, i'_d$ with $k - i'_1 - \cdots - i'_d < m$. To show that $\pi_{i_1,\ldots, i_d} \le m = k - i_1 - \cdots - i_d$ consider two cases. 

{\it Case 1.} If $b_{i_1,\ldots, i_d} \le 1$, then we have 
$$
\pi_{i_1,\ldots, i_d} = b_{i_1,\ldots, i_d} + \max_{\ell \in [d]} \pi_{(i_1, \ldots, i_d) + \mathbf{e}_{\ell}} \le 1 + m - 1 = m
$$
since $\pi_{(i_1,\ldots, i_d) + \mathbf{e}_{\ell}} \le m - 1$ by  assumption.

{\it Case 2.} If $b_{i_1,\ldots, i_d} = 2$, then by the condition we have $b_{j_1, \ldots, j_d}  = 0$ for all $(j_1,\ldots, j_d) = (i_1,\ldots, i_d) + \mathbf{e}_{\ell}$ and $\ell \in [d]$ and hence 
\begin{align*}
\pi_{i_1,\ldots, i_d} %&= b_{i_1,\ldots, i_d} + \max_{\ell \in [d]} \pi_{(i_1, \ldots, i_d) + \mathbf{e}_{\ell}} \\
&= b_{i_1,\ldots, i_d} + \max_{(j_1,\ldots, j_d) = (i_1,\ldots, i_d) + \mathbf{e}_{\ell}} \left( b_{j_1,\ldots, j_d} + \max_{\ell' \in [d]} \pi_{(j_1,\ldots, j_d) + \mathbf{e_{\ell'}}}  \right) \le 2 + m - 2 = m
\end{align*}
since $\pi_{(j_1,\ldots, j_d) + \mathbf{e}_{\ell'}} \le m - 2$ by assumption.

So we showed that $i_1 + \ldots + i_d + \pi_{i_1,\ldots, i_d} \le k$ which means that $D(\pi) \subseteq \Delta_d(k)$ and $\pi \in {A}_d(k)$, i.e. $\Phi : B_d(k) \to A_d(k)$ as needed.
Since $\Phi$ is injective we have $|B_d(k)| \le |A_d(k)|$ which now implies the inequality. 
\end{proof}

\begin{lemma}\label{lemlb}
We have: 
$$
b_d(k) \ge \left(3 \cdot 2^{3^d - 1} \right)^{\binom{\lfloor (k-1)/3 \rfloor}{d}}. 
%\left(2^{2^d} +  2^{2^d - d - 1} \right)^{\binom{\lfloor (k-1)/2 \rfloor}{d}}. 
$$
\end{lemma}

\begin{proof}
Consider 
hypermatrices $B = (b_{\mathbf{i}})_{\mathbf{i} \in \Delta_{d-1}(k-1)} \in B_d(k)$ %with entries $b_{i_1,\ldots, i_d} \in \{0, 1, 2 \}$ 
satisfying the following property: 
%\begin{itemize}
%\item[(1)] 
for $j_1,\ldots, j_d \in \mathbb{Z}_+$ in each $3 \times \cdots \times 3$ ($d$ times) box 
$$J = \{3j_1-2, 3 j_1 - 1, 3j_1 \} \times \cdots \times \{3j_d-2, 3j_d - 1, 3j_d \} \subset \Delta_{d-1}(k-1),$$ %where $j_1,\ldots, j_d$ are even, 
%we have $b_{i_1,\ldots, i_d} \le 1$ if $(i_1,\ldots, i_d) \ne (2j_1- 1, \ldots, 2 j_d - 1)$; 
%there is at most one $(i_1,\ldots, i_d) \in J$ with $b_{i_1,\ldots, i_d} = 2$; 
%\item[(2)] 
%and 
if $b_{\mathbf{i}} = 2$ for $\mathbf{i} \in J$ %= (2j_1 - 1, \ldots, 2 j_d - 1)$, 
then $\mathbf{i} + \mathbf{e}_{\ell} \in J$  and $b_{\mathbf{i} + \mathbf{e}_{\ell}} =0$ for all $\ell \in [d]$. %adjacent with $(i_1, \ldots, i_d)$ (i.e. when $|i_1 - i'_1| + \cdots + |i_d - i'_d| = 1$). 
%\end{itemize}
(We also set $b_{\mathbf{i}} = 0$ for $\mathbf{i}$ which do not belong to any box $J$.)
Note that for different $j_1,\ldots, j_d \in \mathbb{Z}_+$ the boxes $J$ are disjoint and the number of such boxes $J$ is equal to $\binom{\lfloor (k-1)/3 \rfloor}{d}$ (since it is the number of $j_1,\ldots, j_d$ satisfying $3j_1 + \ldots + 3j_d \le k - 1$). For every such box $J$ there are $2^{3^d}$ sub-hypermatrices $(b_{\mathbf{i}})_{\mathbf{i} \in J}$ with $b_{\mathbf{i}} \in \{0, 1\}$ and $2^{d} \times 2^{3^d - d - 1}$ sub-hypermatrices $(b_{\mathbf{i}})_{\mathbf{i} \in J}$ which contain one entry $b_{\mathbf{i}} = 2$ for $\mathbf{i} \in J$ (since  $\mathbf{i} = (i_1,\ldots, i_d) \in J$ can be chosen in $2^d$ ways where $i_{\ell} \in \{3j_{\ell} - 1, 3 j_{\ell} - 2 \}$, %is unique %in $2^d$ ways 
and the remaining $2^{3^d - d - 1}$ elements of $J \setminus S$ where $S = \{\mathbf{i} \} \cup \{\mathbf{i} + \mathbf{e}_{\ell}  : \ell \in [d]\}$ %non-adjacent to $(i_1,\ldots, i_d)$ 
can be filled in $B$ with  $\{0,1 \}$ arbitrarily according to the described property). Therefore, we have 
$$|B_d(k)| \ge \left(2^{3^d} + 2^{d} \times 2^{3^d - d - 1} \right)^{\binom{\lfloor (k-1)/3 \rfloor}{d}}$$
as needed.
\end{proof}


\begin{remark}
%Just following the argument presented in Case 1 in the proof, i.e. filling the 
Counting in $B_d(k)$ only 
hypermatrices $(b_{\mathbf{i}})_{\mathbf{i} \in \Delta_{d-1}(k-1)}$ with $b_{\mathbf{i}} \in \{0,1\}$ gives a weaker inequality $b_d(k) \ge 2^{|\Delta_{d - 1}(k - 1)|} = 2^{\binom{k - 1}{d}}.$
%$$
%a_d(k) \ge 2^{\binom{k-1}{d}}.
%$$
\end{remark}


\begin{remark}\label{aup}
To get an upper bound on $a_d(k)$, it is not difficult to show that for $k \ge d$ we have
$$
a_{d}(k) \le a_{d-1}(k-1) \cdot a_{d-1}(k-2) \cdots a_{d-1}(d).
$$
From this inequality, by induction on $d$ we obtain that 
\begin{align*}%\label{aup}
a_d(k) \le 2^{2 \binom{k-1}{d}}. %2^{2 \frac{k^d}{d!} - O(k^{d-1})}.
\end{align*}
\end{remark}


\begin{lemma}\label{lempa}
We have: 
$$\widetilde{\mathsf{p}}_{d}(n) \ge a_{d}(k) \text{ for } k  = \lfloor ((d+1)!\, n)^{1/(d+1)} \rfloor.$$ 
\end{lemma}
\begin{proof}
Note that 
$$|\Delta_{d}(k)| \le \frac{k^{d+1}}{(d+1)!} \le n$$
and so every $d$-dimensional partition in the set $A_{d}(k)$ has volume at most $n$, which proves the inequality.
%for $|\Delta_d(k)| \le n < |\Delta_d(k + 1)|$.
\end{proof}

%Let 
%$$
%\delta_d = \liminf_{k \to \infty} \frac{\log a_d(k)}{k^d/d!}.
%$$

\begin{theorem}\label{coral}
We have:
%As $n \to \infty$, 
$$
\liminf_{n \to \infty} \frac{\log \widetilde{\mathsf{p}}_{d}(n)}{n^{d/(d+1)}} \ge 
%{|\Delta_{d-1}(k-1)|} = 
\alpha_d,  %\frac{d!^{(d-1)/d}}{(d-1)!} 
$$
where 
$$
 \qquad \alpha_d =  \frac{(d+1) }{(d+1)!^{1/(d+1)}} \cdot \delta_d, \qquad 
 \delta_d = \liminf_{k \to \infty} \frac{\log a_d(k)}{k^d/d!} \ge \log 2 + 3^{-d} \log \frac{3}{2}. %\frac{1}{2^d} \log(2^{2^d} + 2^{2^d - d - 1}). %O\left(\frac{1}{n^{1/d}} \right)
$$
\end{theorem}
\begin{proof}
For $k  = \lfloor ((d+1)!\, n)^{1/(d+1)} \rfloor$ we have $n^{d/(d+1)} \sim (d+1)!^{-d/(d+1)} k^d = \frac{d!}{(d+1)!^{d/(d+1)}}\cdot k^d/d! = \frac{(d+1)!^{1/(d+1)}}{(d+1)} \cdot k^d/d!$ and hence
\begin{align*}
\liminf_{n \to \infty} \frac
{\log \widetilde{\mathsf{p}}_{d}(n)}
{n^{d/(d+1)}} 
&\ge  \frac{(d+1)}{(d+1)!^{1/(d+1)}} \liminf_{k \to \infty} \frac{\log a_{d}(k)}{k^d/d!} = \frac{(d+1)}{(d+1)!^{1/(d+1)}} \cdot \delta_d = \alpha_d.  
%&\ge 
%%{|\Delta_{d-1}(k-1)|} 
%\binom{\lfloor (k-1)/3 \rfloor }{d} \log \left(2^{3^d} \cdot \frac{3}{2} \right)
%= \delta_d \frac{k^{d}}{d!} - O(k^{d-1}) \\
%&= \frac{(d+1)!^{d/(d+1)}}{d!} \cdot \delta_d\, n^{d/(d+1)} - O(n^{(d-1)/(d+1)})
\end{align*}
Using the above lemmas we have 
$$
\delta_d = \liminf_{k \to \infty} \frac{\log a_d(k)}{k^d/d!} \ge \liminf_{k \to \infty} \frac{\log b_d(k)}{k^d/d!} \ge \lim_{k \to \infty} \frac{\binom{\lfloor (k-1)/3 \rfloor }{d} \log \left(2^{3^d} \cdot \frac{3}{2} \right)}{k^d/d!} = \log 2 + 3^{-d} \log \frac{3}{2}.
$$
%\begin{align*}
%%\frac
%{\log \widetilde{\mathsf{p}}_{d}(n)}
%%{n^{(d-1)/d}} 
%&\ge \log a_{d}(k) \ge \log b_d(k)  \\
%&\ge 
%\binom{\lfloor (k-1)/3 \rfloor }{d} \log \left(2^{3^d} \cdot \frac{3}{2} \right)
%= \delta_d \frac{k^{d}}{d!} - O(k^{d-1}) \\
%&= \frac{(d+1)!^{d/(d+1)}}{d!} \cdot \delta_d\, n^{d/(d+1)} - O(n^{(d-1)/(d+1)})
%\end{align*}
%Since $\alpha_d = \frac{(d+1)!^{d/(d+1)}}{d!} \cdot \delta_d = \frac{(d+1)}{(d+1)!^{1/(d+1)}} \cdot \delta_d$ 
%we get the result.
\end{proof}

\begin{remark}
One can also lower bound $\widetilde{\mathsf{p}}_{d}(n)$ via the number of partitions whose diagram lies inside the box $[m]^{d+1}$ for $m^{d+1} \le n$. Bounds on the number of partitions inside a box were obtained in \cite{ms}. However this gives a lower bound smaller than $\alpha_d$. 
\end{remark}

\begin{remark}
To compare our lower bound constant $\alpha_d$ with the constant $\gamma_d$ from MacMahon's numbers, a direct computation shows that  
$\alpha_7 > 1.47347 > \gamma_7 \approx 1.45831$. It is a routine exercise to show that $\alpha_d > \gamma_d$ for $d \ge 7$. 
\end{remark}

\begin{remark}
The constant $\alpha_d$ can be further slightly improved by giving a better lower bound in Lemma~\ref{lemlb} (e.g. by filling more $2$'s for hypermatrices in $B_d(k)$ or dividing into boxes as $J$ of other sizes), which we did not pursue here.   
\end{remark}

\begin{remark}
The upper bound on $a_d(k)$ mentioned in Remark~\ref{aup} tells that bounding the number of partitions inside the simplex $\Delta_d$ cannot give a lower bound constant on $n^{-d/(d+1)} {\log \mathsf{p}}_d(n)$ better than $\frac{(d+1)}{(d+1)!^{1/(d+1)}} \cdot 2 \log 2$. %$2 \alpha_d$.
\end{remark}

\begin{problem}
Find %Prove that there exists 
$
\lim_{k \to \infty} ({k^d/d!})^{-1} {\log a_d(k)} \in (\log 2, \log 4]
$
and prove its existence. %find it.
\end{problem}


\section{An upper bound via vector partitions}\label{secvec}

Recall that vector partition numbers $\mathsf{p}(n_1,\ldots, n_d)$ are defined via the multivariate generating series 
$$
\sum_{n_1,\ldots, n_d } \mathsf{p}(n_1,\ldots, n_d)\, x_1^{n_1} \cdots x_d^{n_d}  = \prod_{i_1,\ldots, i_d \ge 1} \left(1 - x_1^{i_1} \cdots x_d^{i_d} \right)^{-1}.
$$

Recall also that for a partition $\pi \in \mathcal{P}^{(d)}$, the set of {corners} of $\pi$ is defined as follows:
$$
\mathrm{Cor}(\pi) := \left\{\mathbf{i} \in \mathbb{Z}^{d+1}_+ : \mathbf{i} \in D(\pi),\, \mathbf{i} + \mathbf{e}_{\ell} \not\in D(\pi) \text{ for all } \ell \in [d] \right\}.
$$
%(Here $\{e_{\ell}\}$ is the standard basis in $\mathbb{Z}^{d+1}$.) For example, if $\pi \in \mathcal{P}^{(2)}$ is a plane partition then $\mathrm{Cor}(\pi) = \{(i,j,k) : (i,j, k) \in D(\pi), (i+1, j, k), (i, j+1, k) \not\in D(\pi) \}$. %Let $\mathrm{cor}(\pi) := |\mathrm{Cor}(\pi)|$ be the number of corners of $\pi$. 
%The inverse map $\Phi^{-1}$ can be described via projection of corners as follows. Let $\pi \in \mathcal{P}^{(d)}$ and $A = (a_{i_1,\ldots, i_d}) = \Phi^{-1}(\pi)$. Then we have 
%$$
%a_{i_1,\ldots, i_d} = \left| \left\{i : (i_1,\ldots, i_d, i) \in \mathrm{Cor}(\pi) \right\} \right|.
%$$
%Using this, the following multivariate identity was proved in \cite{ay}. %using the map $\Phi$. 
%\begin{theorem}[\cite{ay}]
%The following multivariate identity holds:
%$$%\begin{align}\label{id1}
%\sum_{\pi \in \mathcal{P}^{(d)}} \prod_{(i_1,\ldots, i_d, i_{d+1}) \in \mathrm{Cor}(\pi)} x_{i_1,\ldots, i_d} = \prod_{i_1,\ldots, i_d \ge 1} \left(1 - x_{i_1,\ldots, i_d} \right)^{-1}. \eqno{(\diamond \diamond)}
%$$%\end{align}
%\end{theorem}
For $\pi \in \mathcal{P}^{(d)}$ and $\ell \in [d]$, define the following statistic: 
$$
c_{\ell}(\pi) := \sum_{(i_1,\ldots, i_{d+1}) \in \mathrm{Cor}(\pi)} i_{\ell}.
$$

We show the following interpretation of vector partition numbers via $d$-dimensional partitions using the c-statistic.

\begin{lemma}\label{wp}
We have: %the following interpretation of vector partition numbers:
$$
\mathsf{p}(n_1,\ldots, n_d) = \left| \left\{\pi \in \mathcal{P}^{(d)} : (c_1(\pi), \ldots, c_d(\pi)) = (n_1,\ldots, n_d) \right\} \right|.
$$
\end{lemma}
\begin{proof}
Firstly, note that by definition of c-statistic we have
$$
\sum_{\pi \in \mathcal{P}^{(d)}} x_1^{c_1(\pi)} \cdots x_d^{c_d(\pi)} = \sum_{\pi \in \mathcal{P}^{(d)}} \prod_{(i_1,\ldots, i_d, i_{d+1}) \in \mathrm{Cor}(\pi)} x_1^{i_1} \cdots x_d^{i_d}.
$$
Now setting $x_{i_1,\ldots, i_d} \to x_1^{i_1} \cdots x_d^{i_d}$ in the identity ($\diamond$) (in \textsection\ref{secphi}) %\eqref{id1} 
we get
\begin{align*}%\label{id2}
\sum_{\pi \in \mathcal{P}^{(d)}} \prod_{(i_1,\ldots, i_d, i_{d+1}) \in \mathrm{Cor}(\pi)} x_1^{i_1} \cdots x_d^{i_d} %&= \sum_{\pi \in \mathcal{P}^{(d)}} x_1^{c_1(\pi)} \cdots x_d^{c_d(\pi)} \\
&=\prod_{i_1,\ldots, i_d \ge 1} \left(1 - x_1^{i_1} \cdots x_d^{i_d} \right)^{-1} \\
&= \sum_{n_1,\ldots, n_d} \mathsf{p}(n_1,\ldots, n_d) x_1^{n_1} \cdots x_{d}^{n_d}.
\end{align*}
Combining these identities we therefore have %Therefore,
$$
\sum_{\pi \in \mathcal{P}^{(d)}} x_1^{c_1(\pi)} \cdots x_d^{c_d(\pi)} = \sum_{n_1,\ldots, n_d} \mathsf{p}(n_1,\ldots, n_d) x_1^{n_1} \cdots x_{d}^{n_d}, 
$$ 
which now implies the result.
\end{proof}

Now we show that the c-statistic does not exceed the volume which will be needed for getting the bounds below.

\begin{lemma}
Let $\pi \in \mathcal{P}^{(d)}$ and $\ell \in [d]$. Then $c_{\ell}(\pi) \le |\pi|$. 
\end{lemma}
\begin{proof}
%Let $\ell \in [d]$ be any fixed direction. 
For each corner $(i_1,\ldots, i_{d+1}) \in \mathrm{Cor}(\pi)$,  consider the set 
$$X_{\ell}(i_1,\ldots, i_{d+1}) :=  \{(i_1, \ldots, i_{\ell - 1}, i, i_{\ell + 1}, \ldots, i_{d+1}) : 1 \le i  \le i_{\ell} \} \subseteq D(\pi).$$ 
Notice that the sets $X_{\ell}(i_1,\ldots, i_{d+1})$ for $(i_1,\ldots, i_{d+1}) \in \mathrm{Cor}(\pi)$ are disjoint. Otherwise, a common element in sets $X_{\ell}$ must be produced from distinct corners with the same coordinates except for $\ell$-th one, which for $\ell \ne d+1$ is impossible by definition of corners. %[Why?] 
Hence, 
$$\bigsqcup_{(i_1,\ldots, i_{d+1}) \in \mathrm{Cor}(\pi)} X_{\ell}(i_1,\ldots, i_{d+1}) \subseteq D(\pi)$$
and we obtain that  
$$
c_{\ell}(\pi) = \sum_{(i_1,\ldots, i_{d+1}) \in \mathrm{Cor}(\pi)} i_{\ell} = \sum_{(i_1,\ldots, i_{d+1}) \in \mathrm{Cor}(\pi)} |X_{\ell}(i_1,\ldots, i_{d+1})| \le |D(\pi)| = |\pi|
$$
as needed. 
\end{proof}

\begin{remark}
Note that this inequality is sharp, e.g. let $\pi$ be the partition whose diagram is the simplex $\Delta_d(k)$ (defined in the previous section). Then we have $c_{\ell}(\pi) = |\pi|$ for all $\ell \in [d]$.  
\end{remark}

Let us also denote
$$
\widetilde{\mathsf{p}}(n_1, \ldots, n_d) := \sum_{i_1 \le n_1, \ldots, i_d \le n_d} \mathsf{p}(i_1,\ldots, i_d) \le n_1 \cdots n_d \cdot \mathsf{p}(n_1,\ldots, n_d). 
$$

Now we are ready to prove our main upper bound. 
\begin{theorem}\label{lempp}
We have the following inequality:
$$
\widetilde{\mathsf{p}}_d(n) \le \widetilde{\mathsf{p}}(\underbrace{n, \ldots, n}_{d \text{ times}}).
$$
\end{theorem}
\begin{proof}
By previous lemma, every $d$-dimensional partition $\pi$ with volume at most $n$ has $c_{\ell}(\pi) \le n$ for all $\ell \in [d]$.
Hence we have 
$$
\left\{\pi \in \mathcal{P}^{(d)} : |\pi| \le n \right\} \subseteq \bigsqcup_{i_1,\ldots, i_d \le n} \left\{\pi \in \mathcal{P}^{(d)} : (c_1(\pi), \ldots, c_d(\pi)) = (i_1,\ldots, i_d) \right\}
$$
Taking cardinalities and using Lemma~\ref{wp} we obtain the inequality.
%every $d$-dimensional partition with volume at most $n$, has 
\end{proof}

The following (asymptotically tight) upper bound on vector partition numbers was proved in \cite{bv} (which is derived  using its generating function). 

\begin{lemma}[\cite{bv}]\label{lembv}
We have the following bound:
$$
\log \mathsf{p}(n_1,\ldots, n_d) \le (d+1) \zeta(d+1)^{1/(d+1)} (n_1 \cdots n_d)^{1/(d+1)}.
$$
\end{lemma}

Theorem~\ref{thm:one} now follows from Theorem~\ref{coral}, Theorem~\ref{lempp} and Lemma~\ref{lembv}. Theorem~\ref{thmc} follows from %Lemma~\ref{lempa} and 
Theorem~\ref{lempp}.

\begin{remark}
Let us see what bounds Theorem~\ref{thm:one} gives for $d  \le 3$.
For $d = 1$ the upper bound is asymptotically tight since the number of integer partitions $\mathsf{p}_1(n)$ satisfies 
$$\frac{\log \mathsf{p}_1(n)}{n^{1/2}} \sim 2 \zeta(2)^{1/2} = \beta_1 = \gamma_1 $$ according to the Hardy--Ramanujan formula. %(where $\zeta$ is the Riemann zeta function). 
For $d = 2$ the number of plane partitions $\mathsf{p}_2(n)$ has the asymptotics 
$$\frac{\log \mathsf{p}_2(n)}{n^{2/3}} \sim \gamma_2 = \frac{3 \zeta(3)^{1/3}}{2^{2/3}}  \approx 2.0094$$ %457 
%(see e.g. \cite{?}) 
while the above bounds range about $(1.2187, 3.1898)$. 
For $d = 3$ the asymptotics of the number $\mathsf{p}_3(n)$ of solid partitions is unknown, and the bounds we obtain give that
$$
1.2797 
< \frac{\log \mathsf{p}_{3}(n)}{n^{3/4}} < 
4.0799
$$
for sufficiently large $n$. It was estimated in \cite{dg} that 
%Experimental estimates in \cite{dg} suggest that 
$n^{-3/4}{\log \mathsf{p}_{3}(n)}{} \sim 1.822 \pm 0.001$.
\end{remark}

\subsection{An upper bound via MacMahon's numbers}\label{secmac}
For $\pi \in \mathcal{P}^{(d)}$, the {\it corner-hook volume} statistic $|\cdot|_{ch} : \mathcal{P}^{(d)} \to \mathbb{N}$ is defined as follows:
$$
|\pi|_{ch} := \sum_{(i_1,\ldots , i_{d+1})\, \in\, \mathrm{Cor}(\pi)} (i_1 +\ldots + i_d - d + 1) = c_{1}(\pi) + \cdots + c_{d}(\pi) - (d-1)\cdot |\mathrm{Cor}(\pi)|.
$$
MacMahon's numbers have the following interpretation via $d$-dimensional partitions obtained in \cite{ay}:
$$
\mathsf{m}_d(n) = \left| \left\{\pi \in \mathcal{P}^{(d)} : |\pi|_{ch} = n \right\} \right|.
$$
Similarly denote 
$
\widetilde{\mathsf{m}}_d(n) := \sum_{i \le n}^{} \mathsf{m}_d(i). %< n\, m_d(n).
$
Then the same proof as above shows that $|\pi|_{ch} \le d\, |\pi| - (d - 1) |\mathrm{Cor}(\pi)| \le |\pi| - (d - 1)$ and 
$$
\widetilde{\mathsf{p}}_d(n) \le \widetilde{\mathsf{m}}_d(dn - d + 1).
$$
In particular, we also have the inequality $\mathsf{p}_d(n) \le (dn - d + 1)\, \mathsf{m}_d(dn - d + 1)$. 



%\section*{Acknowledgements}

\

%\newpage
\begin{thebibliography}{abcdefghij}
\bibitem[AY20]{ay}
A. Amanov and D. Yeliussizov, MacMahon's statistics on higher-dimensional partitions, arXiv:2009.00592

\bibitem[And98]{andrews}
G. E. Andrews, The theory of partitions, %Vol. 2, 
Cambridge University Press, 1998.

\bibitem[ABMM67]{atkin}
A. O. L. Atkin, P. Bratley, I. G. Macdonald, and J. K. S. McKay, Some computations for $m$-dimensional partitions, Proc. Cambridge Philos. Soc. {\bf 63} (1967) 1097--1100.

\bibitem[BGP12]{bgp}
S. Balakrishnan, S. Govindarajan, and N. S. Prabhakar, On the asymptotics of higher-dimensional partitions, J. Phys. A {\bf 45} (2012), 055001.

\bibitem[BV92]{bv}
I. B\'ar\'any and A. Vershik, On the number of convex lattice polytopes, Geometric and Functional Analysis {\bf 2} (1992), 381--393.

\bibitem[BBS13]{bbs}
K. Behrend, J. Bryan, and B. Szendr\H{o}i, Motivic degree zero Donaldson--Thomas invariants, Invent. Math. {\bf 192} (2013), 111--160.

\bibitem[BPA97]{bhatia}
D. P. Bhatia, M. A. Prasad, and D. Arora, Asymptotic results for the number of multidimensional partitions of an integer and directed compact lattice animals, J. Phys. A {\bf 30} (1997), 2281.

\bibitem[DPS+23]{dai}
F. Dai, A. Prymak, A. Shadrin, V. N. Temlyakov, and S. Tikhonov, On the cardinality of lower sets and universal discretization, J. Complexity (2023), Article 101726.

\bibitem[Ekh12]{ekhad}
S. B. Ekhad, The number of $m$-dimensional partitions of eleven and twelve, published electronically at \url{https://sites.math.rutgers.edu/~zeilberg/mamarim/mamarimhtml/mDimPars.html}

\bibitem[DG15]{dg}
N. Destainville and S. Govindarajan, Estimating the asymptotics of solid partitions, J. Stat. Phys. {\bf 158} (2015), 950--967.

\bibitem[Gov]{boltzmann}
S. Govindarajan, The partitions project \url{http://boltzmann.wikidot.com/the-partitions-project}

\bibitem[Gov13]{gov}
S. Govindarajan, Notes on higher-dimensional partitions, J. Combin. Theory Ser. A {\bf 120} (2013), 600--622.

\bibitem[Kra16]{krat}
C. Krattenthaler, Plane partitions in the work of Richard Stanley and his school, The mathematical
legacy of Richard P. Stanley (2016), 246--277.

\bibitem[Knu70]{knuth}
D. Knuth, A note on solid partitions,  Math. Comp. {\bf 24} (1970) 955--961.

\bibitem[Mac16]{macmahon}
P. A. MacMahon, Combinatory Analysis, Cambridge University Press, Vol. 1 and 2, 1916.

\bibitem[MS14]{ms}
G. Moshkovitz and A. Shapira, Ramsey theory, integer partitions and a new proof of the Erd\H{o}s–Szekeres theorem, Adv. Math. {\bf 262} (2014), 1107--1129.

\bibitem[MR03]{mr}
V. Mustonen and R. Rajesh, Numerical estimation of the asymptotic behaviour of solid partitions of an integer, J. Phys. A 36 {\bf 24} (2003), 6651--6659.

\bibitem[Nek20]{nekrasov}
N. Nekrasov, Magnificent four, Annales de l’Institut Henri Poincar\'e D {\bf 7} (2020), 505--534.

\bibitem[OEIS]{oeis}
Online Encyclopedia of Integer Sequences \url{https://oeis.org}

\bibitem[Sta71]{sta2}
R. P. Stanley, Theory and application of plane partitions, Parts 1 and 2, Studies in Applied Math. {\bf 50} (1971), 167--188, 259--279.

\bibitem[Sta99]{sta}
R. P. Stanley, Enumerative Combinatorics, Vol. 2, Cambridge, 1999.

\end{thebibliography}

\end{document}






Now we show that our lower bound constant $\alpha_d$ is eventually larger than the constant $\gamma_d$ from MacMahon's numbers. 

\begin{lemma}\label{lemag}
$\alpha_d > \gamma_d$ for $d \ge 7$.
%Let $d \ge 7$. We have 
%$$
%\frac{(d+1)}{(d+1)!^{1/(d+1)}} \log 2 > \frac{d+1}{d} \left(d \zeta(d+1) \right)^{1/(d+1)}
%$$
\end{lemma}
\begin{proof}
A direct computation shows that 
$\alpha_7 > 1.47 > \gamma_7 \approx 1.4583108.$
Let us show that the sequence $f_d = \gamma_d/\alpha_d$ is decreasing for positive integers $d$, which will then imply the needed inequality since we get that $f_d < f_7 < 1$ for $d > 7$. 
We have
$$\frac{f_d}{f_{d+1}} = \frac{\gamma_d}{\gamma_{d+1}} \cdot \frac{\alpha_{d+1}}{\alpha_d}= \frac{d^{-d/(d+1)} (d+1) \zeta(d+1)^{1/(d+1)}}{ (d+1)^{-(d+1)/(d+2)} (d+2) \zeta(d+2)^{1/(d+2)}} \cdot \frac{  (d+2)!^{-1/(d+2)} (d+2)^{}  }{   (d+1)!^{-1/(d+1)} (d+1)^{} }$$
Since $\zeta(d+1) > \zeta(d+2) > 1$ we have $\zeta(d+1)^{1/(d+1)} > \zeta(d+2)^{1/(d+2)}$. Using this and simplifying we have
$$
\frac{f_d}{f_{d+1}} > \frac{(d+1)^{(d+1)/(d+2)}}{d^{d/(d+1)}} \cdot \frac{(d+1)!^{\frac{1}{(d+1)(d+2)} }}{ (d+2)^{1/(d+2)} }
$$
Using the known bound %on factorial 
$
(d+1)! \ge \frac{(d+1)^{(d+1)}}{e^d}  %\le \frac{(d+1)^{(d+2)}}{e^d}.
$
we obtain that
$$
\frac{f_d}{f_{d+1}} > \frac{(d+1)^{(d+1)/(d+2)}}{d^{d/(d+1)}} \cdot \frac{(d+1)^{\frac{1}{(d+2)} }}{ (d+2)^{1/(d+2)} e^{\frac{d}{(d+1)(d+2)}} } = \frac{(d+1)}{d^{d/(d+1)} (d+2)^{1/(d+2)} e^{\frac{d}{(d+1)(d+2)}} }.
$$
We also have the bound $e^{-1/(d+1)} \ge  \frac{d}{d +1 }$ using which we obtain that
$$
\frac{f_d}{f_{d+1}} > \frac{(d+1) \cdot d^{d/(d+2)}}{d^{d/(d+1)} (d+2)^{1/(d+2)} (d+1)^{d/(d+2)} } 
= \frac{(d+1)^{2/(d+2)}  }{d^{\frac{d}{(d+1)(d+2)}} (d+2)^{1/(d+2)} } %> \frac{d(d+2)^{1/(d+2)}  }{d^{\frac{d}{(d+1)(d+2)}} (d+2)^{1/(d+2)} }
$$
Since $(d+1)^2 > d(d+2)$ we get
$$
\frac{f_d}{f_{d+1}} >\frac{(d(d+2))^{1/(d+2)}  }{d^{\frac{d}{(d+1)(d+2)}} (d+2)^{1/(d+2)} } = d^{\frac{1}{(d+1)(d+2)}} \ge 1
$$
which gives the needed. 
\end{proof}


